\documentclass[11pt]{article}
\usepackage{amsmath,amssymb,geometry}
\geometry{margin=22mm}
\usepackage[UTF8]{ctex}
\begin{document}
\section*{P56：选择了 Euler 连续松弛算法，并定义了 CVaR 值。}
该模型保留了最初选定的 M1–M5 和 M8 参数，包括其明确的一阶欧拉状态转移以及连续控制域。风险则通过源论文中的方程（6）中的有限情景 CVaR 来表示。

\subsection*{数据和变量}
有$J$个控制通道、$T$个时间步长，以及$S$种情景；还有$\Delta t>0$、$p_s\ge0$、$\sum_s p_s=1$、$\eta\in(0,1]$、$\alpha\in[0,1]$和$\theta\ge0$。  
漂移/控制矩阵为$H^{(0)},H^{(j)}\in\mathbb C^{d\times d}$，初始矩阵为$X_{\rm init}$，目标值为$X_{\rm targ}$，而实际扰动则表示为$\xi_{jk}^s$。  
选择$u_{jk}\in[0,1]$、复态$X_k^s$、实数阈值$\zeta$，以及非负数值$r_s$。所有矩阵维度和情景/时间索引均被保留下来。

\subsection*{完整的程序如下：}
\begin{align}
F_L(u)&=\sum_{k=1}^T\left(\sum_{j=1}^J u_{jk}-1\right)^2, \tag{M1}\\
H_k^s&=(1+\xi_{0k}^s)H^{(0)}
       +\sum_{j=1}^J(1+\xi_{jk}^s)u_{jk}H^{(j)}, \tag{M2}\\
X_k^s&=(I-\mathrm i\Delta t H_k^s)X_{k-1}^s
      &&(s=1,\ldots,S;\ k=1,\ldots,T), \tag{M3}\\
X_0^s&=X_{\rm init} &&(s=1,\ldots,S), \tag{M4}\\
F_s&=1-\frac{|\operatorname{tr}(X_{\rm targ}^H X_T^s)|}{d}
      &&(s=1,\ldots,S), \tag{M5}\\
r_s&\ge F_s-\zeta,\qquad r_s\ge0 &&(s=1,\ldots,S), \tag{M6R}\\
F_C&=\zeta+\eta^{-1}\sum_s p_sr_s. \tag{M7R}
\end{align}
最小化
\[
 \alpha\sum_s p_sF_s+(1-\alpha)F_C+\theta F_L(u)
 \quad\text{subject to M1--M6R and }u\in[0,1]^{J\times T},\
 \zeta\in\mathbb R. \tag{M8}
\]
该轨迹表达式用源方程（3）完成了之前符号化的$F_X$。对于任何一个固定的$u$，最优的CVaR值等于……
\[
 \min_{z\in\{F_1,\ldots,F_S\}}
 \left(z+\eta^{-1}\sum_s p_s\max\{F_s-z,0\}\right).
\]
求解器通过 M3 消除了这些状态，同时通过此断点减少了 CVaR 变量的值。不过，这两种消除方式都没有对修复后的程序产生任何影响。

\subsection*{近似值和求解器限制}
该模型通过一阶欧拉递推公式来近似描述量子演化过程，同时采用连续控制方式。基于迹的损耗值为实数，且可能为负数。通道和偏差会引入平方形式的惩罚项。状态的可行性则由递推公式以及所显示的变量域来决定。

六次可重复进行的有限域局部搜索，能够在 3600 秒的预算时间内返回最佳可行点；收敛过程可能会提前结束。求解器会报告最佳可行的局部搜索点。另外，一个独立的实/虚递归过程会独立检查所有状态、控制边界、损失、CVaR 图/值、惩罚以及目标函数。单位偏差会被记录为诊断信息。
\end{document}
